\section{ESPEON}
Engajamento e Supervisão do Processo de Ensino Online (ESPEON) é uma extensão para navegadores Chrome, em desenvolvimento para medir os níveis de atenção e engajamento de alunos em aulas remotas. Será utilizada para validar a efetividade da metodologia expositiva no contexto do ensino remoto.

\section{Metodologia Expositiva}
\citet{andreataAulaExpositivaPaulo2019} define a metodologia expositiva como aquela baseada na transmissão do saber por parte de um orador, por meio da exposição do conteúdo, enquanto ao estudante cabe uma postura predominantemente passiva de recepção e assimilação do conhecimento. 

\section{MongoDB}
MongoDB é um banco de dados não relacional de propósito geral, orientado a documentos, amplamente utilizado no desenvolvimento de aplicações modernas e escaláveis, especialmente em ambientes de computação em nuvem.

\section{AtlasDB}
AtlasDB é um serviço de banco de dados na nuvem oferecido pela MongoDB Inc., que facilita a implantação, o gerenciamento e o escalonamento de instâncias MongoDB.

\section{API}
Sigla de \textit{Application Programming Interface}, uma API é um conjunto de definições e protocolos que permite a comunicação entre sistemas, aplicativos ou serviços.

\section{Flask}
Flask é um framework da linguagem Python voltado para aplicações web. Ele permite o desenvolvimento ágil de sistemas robustos, com suporte à escalabilidade e integração com diversas bibliotecas.

\section{\textit{Deploy}}
Termo em inglês comumente traduzido como ``implantação'', refere-se ao processo de disponibilização de um sistema ou aplicação em ambiente de produção.

\section{\textit{Log}}
Em computação, \textit{logs} são registros estruturados que documentam eventos ou ações executadas por um sistema, comumente utilizados para auditoria, monitoramento e análise de comportamento.

% \section{\textit{Payload}}


\section{Chrome Web Store}
Plataforma oficial de distribuição de extensões, aplicações e temas para o navegador Google Chrome.

